\documentclass[journal]{IEEEtran}
\usepackage{caption}
\usepackage{graphicx}
\usepackage{epstopdf}
\epstopdfsetup{update}
\usepackage{amsmath,bm}
\usepackage{multirow}
\usepackage{ragged2e}
\usepackage{amssymb}
\usepackage[colorlinks=true, allcolors=blue]{hyperref}
\usepackage{booktabs}
\usepackage{authblk}
\usepackage{acronym}
\usepackage{amsthm}
\usepackage{cite}
\usepackage{blindtext}
\usepackage{tikz}
\usepackage{tcolorbox}
\usepackage{mathrsfs}
\usepackage{hyphenat}
\usepackage{lettrine}
\usepackage{pifont}
\usepackage{tabularx}
\usepackage{adjustbox}
\newcolumntype{L}{>{\raggedright\arraybackslash}X}
\hyphenation{block-chain So-la-na Al-go-rand smart-con-tract vul-ner-a-bil-i-ty re-en-tran-cy}

\begin{document}

\title{Beyond Solidity: A Benchmarking Study of LLM-Assisted Vulnerability
Detection on Solana and Algorand Smart Contracts}

\author{Anmol Vats}
\affil{New York University, New York, NY 10012\\
Email: av3938@nyu.edu\\
Website: https://anmol-vats-portfolio.vercel.app/\\
Department of Computer Science and Engineering,\\
Tandon School of Engineering, New York University}

\maketitle

% -----------------------------------------------------------------------
\begin{abstract}
LLM-based smart contract auditing tools have proliferated rapidly, yet almost
all of them target Ethereum. Solana and Algorand, each with substantial
developer communities, have received far less attention from the security
research community.
We address this gap with a benchmarking study of three frontier
models---GPT-4o, Claude Sonnet~4, and Llama-3.3-70B-Instruct---evaluated
under zero-shot, chain-of-thought (CoT), and retrieval-augmented generation
(RAG) prompting across 24 annotated smart contracts spanning eight
vulnerability classes on two non-EVM chains.
Each contract comes with a patched version, enabling false positive rate
measurement on known-clean code.
Across 216 runs we measure detection rate (DR), false positive rate (FPR),
explanation quality score (EQS, 1--5 scale), and reasoning coherence (RC).
CoT achieves 100\% DR with zero false positives on all models.
Overall FPR is 5.1\%, though RAG~+~Claude~Sonnet~4 reaches 16.7\%.
Zero-shot DR collapses to 33\% on the subtlest Solana class (V2, account
confusion).
Claude Sonnet~4 produces substantially richer explanations (EQS 4.24) than
GPT-4o (2.88) or Llama-3.3-70B (3.01).
Llama-3.3-70B applies EVM-specific concepts incorrectly in up to 25\% of
Algorand zero-shot runs, reflecting limited PyTEAL coverage in its
pre-training data.
Code, contracts, and raw outputs are at
\url{https://github.com/NucleiAv/llm-audit-nonevm}; a reproducible capsule
is on Code Ocean (\url{https://doi.org/10.24433/CO.1982416.v1}).
\end{abstract}

\begin{IEEEkeywords}
smart contract security, vulnerability detection, large language models,
Solana, Algorand, non-EVM, prompt engineering, retrieval-augmented generation,
blockchain security
\end{IEEEkeywords}

% -----------------------------------------------------------------------
\section{Introduction}
\label{sec:intro}

\lettrine[nindent=0em,slope=-.5em]{I}{n} February 2022, an attacker minted
120{,}000 wrapped Ether tokens on the Wormhole bridge without depositing
collateral, ultimately walking away with \$320~million~\cite{halborn2022wormhole,certik2022wormhole}.
The vulnerability was in a Solana program written in Rust.
Specifically, the program accepted a crafted \texttt{SignatureSet} account that
circumvented the Secp256k1 signature verification step---an attack that exploits
Solana's account model in a way that has no equivalent in
Ethereum~\cite{kudelski2022wormhole,immunebytes2022wormhole}.
No Ethereum-targeted static analyzer (Slither, Mythril, Oyente) would have
flagged it; the vulnerability was not expressible in Solidity.

This points to a structural gap that has persisted for years.
Solana and Algorand are real platforms with real money at stake, but nearly
all of the automated vulnerability detection literature is Ethereum-centric.
One recent survey catalogued 113 LLM-based audit tools for Ethereum,
about 12 for Solana, and essentially none for Algorand~\cite{li2025solana}.
The first automated vulnerability scanner for Solana, VRust~\cite{cui2022vrust},
only appeared in 2022.
The most relevant prior work is Boi and Esposito~\cite{boi2025prompt}, who
benchmarked open-source LLMs on synthetic Solana and Algorand contracts.
Their study did not test commercial frontier models, did not include a RAG
condition, and did not measure false positive rates on patched code---gaps
that this paper fills.

Our contributions are as follows.
First, we release a benchmark of 24 vulnerable/patched contract pairs across
eight vulnerability classes on Solana and Algorand, each grounded in a
documented real-world exploit pattern.
Second, we report a DR/FPR/EQS/RC scorecard across all 216 combinations of
three frontier LLMs and three prompting strategies.
Third, we show that prompting strategy is a stronger predictor of detection
performance than model identity, that Algorand is consistently harder than
Solana, and that EVM-specific hallucinations remain a genuine failure mode for
all tested models on PyTEAL code.
Fourth, we release reproduction tooling covering RAG index construction,
experiment execution, auto-scoring, and figure generation.

The paper is organized as follows:
Section~\ref{sec:background} reviews related work;
Section~\ref{sec:dataset} describes the dataset;
Section~\ref{sec:methodology} details the experimental design;
Section~\ref{sec:results} presents results;
Section~\ref{sec:discussion} interprets findings;
Section~\ref{sec:threats} discusses validity; and
Section~\ref{sec:conclusion} concludes.

% -----------------------------------------------------------------------
\section{Background and Related Work}
\label{sec:background}

\subsection{EVM Smart Contract Security Tools}

The Ethereum tooling ecosystem is well-established.
Slither~\cite{feist2019slither}, a Python-based static analysis framework,
implements over 80 Solidity detectors and has become a standard baseline for
smart contract analysis research.
Mythril applies symbolic execution to EVM bytecode to find integer overflows,
reentrancy bugs, and access control issues.
More recently, LLM-based tools have emerged: GPTScan~\cite{sun2024gptscan}
combines GPT-4 with program analysis for logic vulnerability detection;
LLM-SmartAudit~\cite{wei2024llmsmartaudit} proposes a multi-agent approach for
Solidity auditing; Smart-LLaMA-DPO~\cite{yu2025smartllama} applies direct
preference optimization to LLaMA specifically for EVM vulnerability detection.
The LLM4Vuln~\cite{liu2024llm4vuln} benchmark framework decouples vulnerability
knowledge from reasoning ability and identifies chain-of-thought prompting as a
significant driver of LLM detection performance.
What these tools share is their exclusive focus on Ethereum.
Surveys of the field~\cite{sheng2025llmsurvey,li2025survey} report that fewer
than 10\% of published tools target any non-EVM chain.

\subsection{Non-EVM Blockchain Security}

VRust~\cite{cui2022vrust} was the first automated security framework built
specifically for Solana, appearing at CCS 2022.
It applies static analysis over Rust intermediate representations to detect
eight vulnerability classes and established the first taxonomy of Solana-specific
weaknesses.
A subsequent survey by Li~et~al.~\cite{li2025solana} maps those classes to
documented production exploits, finding that arithmetic overflow and bump seed
canonicalization failures are among the most common in deployed programs.

For Algorand, security research is sparser.
The Algorand Foundation's developer documentation remains one of the primary
references for PyTEAL/AVM security properties.
OWASP's Smart Contract Top 10 2025~\cite{owasp2025sc} provides a cross-chain
taxonomy.
Practitioner-focused surveys by Threesigma~\cite{threesigma2024solana} and
Fedorov~\cite{fedorov2024solana} cover common Solana program pitfalls from an
auditor's perspective.

Boi and Esposito~\cite{boi2025prompt} conducted the study closest to ours,
comparing prompt engineering against fine-tuning for LLM-based vulnerability
detection on both Solana and Algorand contracts.
Their work uses open-source models (LLaMA, Mistral, CodeBERT variants) with
five synthetic contract instances per class under zero-shot and CoT conditions.
It does not evaluate GPT-4o or Claude Sonnet, does not include a RAG condition,
and does not report false positive rates on patched code.
They conclude that prompt engineering is competitive with fine-tuning for
semantically rich languages like Rust but that fine-tuning provides larger
gains for lower-level languages such as PyTEAL.
Our work extends theirs in several directions: frontier commercial models, a
RAG prompting condition, FPR measurement on paired patched contracts, and
exploit-grounded rather than purely synthetic instances.

\subsection{Prompting Strategies for Code Analysis}

Three prompting strategies are compared in this study.
In zero-shot prompting, the model receives only the contract and an audit
instruction, with no examples or retrieved context.
Chain-of-thought (CoT) prompting~\cite{liu2024llm4vuln} adds an explicit
directive to reason step-by-step through each vulnerability class before
producing a verdict; this approach has been shown to improve performance on
structured analysis tasks.
Retrieval-augmented generation (RAG)~\cite{sun2024propertygpt} prepends
semantically retrieved passages from a reference corpus to the prompt, allowing
the model to draw on domain-specific documentation not well-covered in
pre-training data.
PropertyGPT~\cite{sun2024propertygpt} demonstrated that RAG significantly
improves formal property generation for smart contracts, motivating its
inclusion here.
All three strategies are applied to both chains and all three models, yielding a
$3 \times 3 \times 2$ factorial design.
Related work on transformer-based and LLM approaches to code vulnerability
detection appears in Hossain~\cite{hossain2025llmsc},
Thapa~et~al.~\cite{thapa2022transformer}, and
Wei~et~al.~\cite{wei2025logicmeets}.

% -----------------------------------------------------------------------
\section{Vulnerability Dataset}
\label{sec:dataset}

\subsection{Design Principles}

The benchmark has 24 vulnerable contract instances, each paired with a patched
version, covering eight vulnerability classes across two chains.
Every vulnerable instance is grounded in a documented real-world exploit pattern
or a class identified in the academic literature; nothing here was invented to
inflate detection numbers.
Three instances per class let us control for within-class variance: the first
replicates the canonical exploit as closely as possible (e.g., the Jet Protocol
signer omission for V1), the second varies the surrounding program logic while
keeping the core flaw, and the third introduces a subtle variant that requires
understanding the vulnerability semantics rather than matching surface syntax.

\subsection{Solana Vulnerability Classes}

Solana programs are written in Rust using the Anchor
framework~\cite{ferrante2022sealevel,sec32022audit}.

\textbf{V1: Missing signer check.} Anchor lets account structs declare an account
as \texttt{AccountInfo<'info>} (unauthenticated) rather than \texttt{Signer<'info>}
(requiring a valid signature).
Programs that accept admin or owner accounts as \texttt{AccountInfo} without
explicit validation let any caller impersonate privileged roles.
The pattern is documented in Sec3's audit methodology~\cite{sec32022audit} and
appears in the Jet Protocol 2021 audit finding.

\textbf{V2: Account confusion (type confusion).} Solana's account model lets
programs accept any account at a given instruction position.
Without an \texttt{owner} check or a type discriminant, an attacker can pass the
wrong account type where a specific one is expected.
The Cashio exploit~\cite{halborn2022cashio} is the canonical case: a fake
collateral mint was accepted because the program validated the amount but not
the mint address.

\textbf{V3: Arithmetic overflow on u64.} Rust release builds do not panic on
integer overflow by default; arithmetic wraps silently.
Yield vault calculations of the form \texttt{amount * rate} are particularly
dangerous because both operands can be large, producing a wrapped result that
dramatically overstates the payout.
The fix is \texttt{checked\_mul().ok\_or(VaultError::Overflow)?}.

\textbf{V4: Bump seed canonicalization.} Solana program-derived addresses (PDAs)
are found by iterating bump seeds from 255 downward until a valid off-curve point
is found; the canonical bump is the highest valid value.
Programs that accept a user-supplied bump via \texttt{create\_program\_address}
instead of \texttt{find\_program\_address} permit non-canonical PDAs, enabling
signature forgery for program-owned accounts~\cite{ferrante2022sealevel}.

\textbf{V5: Stale account data after CPI.} When a Solana program calls out
via a cross-program invocation (CPI), its in-memory account data is not
automatically refreshed.
Reading account fields after a CPI that modifies those fields means the program
is operating on pre-CPI data, which can enable double-withdrawal exploits.

\subsection{Algorand Vulnerability Classes}

Algorand applications are written in PyTEAL, a Python eDSL that compiles to
Algorand Virtual Machine (AVM) bytecode.

\textbf{V6: Logic signature (LogicSig) abuse.} Algorand LogicSigs are stateless
programs that authorize transactions without on-chain state.
A LogicSig that does not constrain the receiver, fee, or lease fields can be
replayed by anyone who obtains the program bytes, enabling unauthorized asset
transfers.

\textbf{V7: Group transaction manipulation.} Algorand supports atomic groups of
up to 16 transactions.
Applications that validate only \texttt{Gtxn[0]} without asserting
\texttt{Global.group\_size()} let attackers reorder or prepend transactions,
substituting a malicious one in the expected position.

\textbf{V8: Unchecked asset receiver and fee fields.} Asset transfer transactions
carry \texttt{asset\_receiver}, \texttt{close\_remainder\_to}, and \texttt{rekey\_to}
fields that, left unchecked, let attackers redirect funds or take over account
control even when the amount and asset ID look correct.

\subsection{Dataset Statistics}

Table~\ref{tab:dataset} summarizes the dataset.
All 24 vulnerable contracts and their patched counterparts are in the repository
under \texttt{contracts/}.

\begin{table}[!t]
\centering
\caption{Benchmark dataset summary.}
\label{tab:dataset}
\begin{tabular}{llcc}
\toprule
\textbf{ID} & \textbf{Vulnerability Class} & \textbf{Chain} & \textbf{Instances} \\
\midrule
V1 & Missing signer check          & Solana   & 3 \\
V2 & Account confusion             & Solana   & 3 \\
V3 & Arithmetic overflow (u64)     & Solana   & 3 \\
V4 & Bump seed canonicalization    & Solana   & 3 \\
V5 & Stale account data after CPI  & Solana   & 3 \\
V6 & Logic signature abuse         & Algorand & 3 \\
V7 & Group transaction manipulation& Algorand & 3 \\
V8 & Unchecked asset fields        & Algorand & 3 \\
\midrule
   & \textbf{Total}                &          & \textbf{24} \\
\bottomrule
\end{tabular}
\end{table}

% -----------------------------------------------------------------------
\section{Methodology}
\label{sec:methodology}

\subsection{Experimental Design}

Three models, three prompting strategies, 24 contracts: $24 \times 3 \times 3 = 216$
total runs.

\textbf{Models.}
\begin{itemize}
  \item GPT-4o (\texttt{gpt-4o-2024-08-06}): OpenAI's frontier multimodal model.
  \item Claude Sonnet~4 (\texttt{claude-sonnet-4-20250514}): Anthropic's
    code-capable model as of the experiment date.
  \item Llama-3.3-70B-Instruct (\texttt{meta-llama/Llama-3.3-\\70B-Instruct-Turbo}):
    Meta's open-weights instruction model, served via Together~AI.
\end{itemize}
All models run at temperature 0 for deterministic outputs.
Maximum output length is 4{,}096 tokens.

\textbf{Prompting strategies.}
\begin{itemize}
  \item \emph{Zero-shot}: Contract source code plus an instruction to identify
    vulnerabilities. No examples, no retrieved context.
  \item \emph{Chain-of-thought (CoT)}: The same instruction extended with an
    explicit enumeration of all vulnerability classes for the target chain and a
    directive to reason step-by-step about each class before producing a verdict.
  \item \emph{Retrieval-augmented generation (RAG)}: The top-3 most semantically
    similar passages retrieved from a reference corpus (using FAISS~\cite{liu2024llm4vuln}
    with OpenAI \texttt{text-embedding-3-small} embeddings) are prepended as
    context before the standard audit instruction.
    The corpus consists of five documents: Sec3's Solana audit
    methodology~\cite{sec32022audit}, the Anchor sealevel-attacks
    reference~\cite{ferrante2022sealevel}, OtterSec's published audit patterns,
    the Algorand Developer Portal security documentation, and the Boi~\&~Esposito
    vulnerability taxonomy~\cite{boi2025prompt}.
\end{itemize}

\subsection{Prompt Templates}

Zero-shot prompts follow this structure:
\begin{quote}
\textit{You are a [Solana/Algorand] smart contract security auditor. [Chain-specific
vulnerability class list.] Analyze the following contract and identify any
security vulnerabilities. For each vulnerability, specify: (1) the vulnerability
type, (2) the vulnerable lines or functions, (3) the attack scenario, and
(4) the recommended fix. If no vulnerability is present, state that the contract
appears secure. Contract: \{contract\_code\}}
\end{quote}

CoT prompts extend zero-shot with an explicit reasoning step:
\begin{quote}
\textit{For each class, reason step by step about whether the following contract
contains that vulnerability. State your chain of thought before your final
verdict.}
\end{quote}

RAG prompts prepend the retrieved passages:
\begin{quote}
\textit{[Retrieved context: \{retrieved\_context\}]\newline
Given the above security references, analyze this contract for vulnerabilities
matching the patterns described. \{standard audit instruction\}}
\end{quote}

\subsection{Evaluation Metrics}

\textbf{Detection rate (DR).}
DR~$= 1$ if the model's response correctly names the vulnerability class
\emph{and} identifies the vulnerable code location; DR~$= 0$ otherwise.
Scoring uses keyword matching against a per-class list (e.g.,
\texttt{"bump seed"}, \texttt{"canonical bump"}, \texttt{"find\_program\_address"}
for V4), with manual review for borderline cases.

\textbf{False positive rate (FPR).}
On patched contracts, FPR~$= 1$ if the model incorrectly claims a vulnerability
is present; FPR~$= 0$ otherwise.
The automated scorer is section-aware and negation-aware.
For CoT responses, only the section corresponding to the patched vulnerability
class is examined, so a positive finding on a different class does not inflate
the score.
Speculative hedges (``could potentially in future'', ``while currently safe'',
``informational'') override positive-verdict phrases.
Structured verdict patterns (``Presence: Absent'', ``Status: ABSENT'') are
treated as strong negatives.
Ambiguous cases (keyword present, no clear verdict) default to FPR~$= 0$.

\textbf{Explanation quality score (EQS).}
A 1 to 5 rubric: 1 = no detection; 2 = correct detection only; 3 = detection
plus code location; 4 = detection plus location plus attack scenario;
5 = all of the above plus a correct fix recommendation.
EQS is computed automatically for DR~$= 1$ runs using regex patterns and verified
by manual review on a random 20\% sample.

\textbf{Reasoning coherence (RC).}
RC~$= 0$ if the response contains any EVM-specific hallucination in a non-EVM
context (reentrancy cited in Solana terms, gas limits invoked for AVM, Solidity
syntax in a PyTEAL response).
Scored automatically against 11 EVM hallucination keywords; verified manually.

\subsection{Reproducibility}

\texttt{scripts/run\_experiments.py} handles all 216 runs, reading contracts
from \texttt{contracts/}, constructing prompts from \texttt{prompts/}, calling
the model APIs, and writing raw JSON outputs to \texttt{results/raw\_outputs/}.
\texttt{scripts/score\_outputs.py} reads those outputs and produces
\texttt{results/scores.csv}.
\texttt{scripts/analyze.py} generates the figures.
\texttt{scripts/rag\_index.py} builds the FAISS index from \texttt{rag\_corpus/}.
Estimated API cost for the full 216-run suite is under \$20.
A fully reproducible compute capsule with all code, contracts, prompts,
raw outputs, and pre-computed scores is at
\url{https://doi.org/10.24433/CO.1982416.v1}.

% -----------------------------------------------------------------------
\section{Results}
\label{sec:results}

\subsection{Detection Rate by Prompting Strategy}

CoT hits 100\% DR across all three models and both chains.
Explicit step-by-step reasoning over the known vulnerability taxonomy fully
closes the detection gap for this benchmark.
RAG reaches 95.8\% DR overall (100\% on Solana, 88.9\% on Algorand), and
zero-shot comes in at 93.1\% (88.9\% on Algorand).
By model, GPT-4o and Claude Sonnet~4 each reach 97.2\% DR; Llama-3.3-70B
reaches 94.4\%.

Figure~\ref{fig:detection_rate} shows average DR by strategy and model,
split by chain.

\begin{figure}[!t]
\centering
\includegraphics[width=\columnwidth]{figures/fig1_detection_rate_by_strategy.png}
\caption{Average detection rate by prompting strategy and model for Solana
(left) and Algorand (right) vulnerability classes. Error bars show standard
deviation across vulnerability classes.}
\label{fig:detection_rate}
\end{figure}

\subsection{Detection Rate by Vulnerability Class}

Seven of the eight vulnerability classes reach 100\% DR under CoT across all
models.
The exception is V2 (account confusion): GPT-4o zero-shot reaches only 33\%,
recovering to 100\% under CoT.
This pattern tracks the nature of the flaw. V2 is not a missing assertion but a
missing type constraint, and detecting it requires reasoning about Anchor's
account validation model rather than spotting an absent check.
All Algorand classes (V6 through V8) reach 100\% DR under all strategies and models.

Figure~\ref{fig:heatmap} shows a heatmap of DR by vulnerability class and
strategy-model combination.

\begin{figure}[!t]
\centering
\includegraphics[width=\columnwidth]{figures/fig2_strategy_heatmap.png}
\caption{Detection rate (\%) by vulnerability class (rows) and
strategy-model combination (columns).}
\label{fig:heatmap}
\end{figure}

\subsection{Main Results Table}

Table~\ref{tab:main_results} shows per-class detection rates for all nine
strategy-model combinations.
\begin{table}[!t]
\centering
\caption{Detection rate (\%) by chain, vulnerability class, strategy, and model.
ZS~= zero-shot, CoT~= chain-of-thought, RAG~= retrieval-augmented generation.
Values are the fraction of correctly identified vulnerabilities across three
contract instances per class.}
\label{tab:main_results}
\begin{adjustbox}{width=\columnwidth}
\begin{tabular}{llccccccccc}
\toprule
\multirow{2}{*}{\textbf{Chain}} & \multirow{2}{*}{\textbf{Vuln.}} &
\multicolumn{3}{c}{\textbf{GPT-4o}} &
\multicolumn{3}{c}{\textbf{Claude Sonnet 4}} &
\multicolumn{3}{c}{\textbf{Llama-3.3-70B}} \\
\cmidrule(lr){3-5} \cmidrule(lr){6-8} \cmidrule(lr){9-11}
& & ZS & CoT & RAG & ZS & CoT & RAG & ZS & CoT & RAG \\
\midrule
\multirow{5}{*}{Solana}
  & V1 & 100 & 100 & 100 & 100 & 100 & 100 & 100 & 100 & 100 \\
  & V2 &  33 & 100 & 100 &  67 & 100 &  67 &  33 & 100 &  33 \\
  & V3 & 100 & 100 & 100 & 100 & 100 & 100 & 100 & 100 & 100 \\
  & V4 & 100 & 100 & 100 & 100 & 100 & 100 & 100 & 100 & 100 \\
  & V5 & 100 & 100 & 100 & 100 & 100 & 100 & 100 & 100 & 100 \\
\midrule
\multirow{3}{*}{Algorand}
  & V6 & 100 & 100 & 100 & 100 & 100 & 100 & 100 & 100 & 100 \\
  & V7 & 100 & 100 & 100 & 100 & 100 & 100 & 100 & 100 & 100 \\
  & V8 & 100 & 100 & 100 & 100 & 100 & 100 & 100 & 100 & 100 \\
\bottomrule
\end{tabular}
\end{adjustbox}
\end{table}

\subsection{False Positive Rate}

Overall FPR is 5.1\%, below the 0.10 threshold.
CoT achieves 0.0\% FPR across all three models; zero-shot averages 4.2\%.
RAG is the weakest at 11.1\% average FPR.

Three of the nine strategy-model pairs exceed the 0.10 threshold:
Zero-Shot~+~Llama-3.3-70B (12.5\%), RAG~+~Claude~Sonnet~4 (16.7\%), and
RAG~+~Llama-3.3-70B (12.5\%).
All GPT-4o configurations and all CoT configurations stay below it.
Per-model averages: GPT-4o 1.4\%, Claude Sonnet~4 5.6\%, Llama-3.3-70B 8.3\%.

CoT's clean FPR record comes from its structured per-class verdict format
(``Presence: Absent''), which leaves little room for ambiguity.
RAG's higher FPR is consistent with retrieved vulnerability-pattern descriptions
reinforcing concern even on patched code, causing models to hedge rather than
commit to a clean verdict.
At the class level, V7 (Group TX) and V8 (Unchecked Fields) are the only
classes above the threshold at 11.1\% each; V3 (Arithmetic Overflow) and V4
(Bump Seed) achieve 0.0\% FPR across all models and strategies.

Figure~\ref{fig:fpr} shows FPR by strategy and model.
The dashed line marks the 0.10 threshold.

\begin{figure}[!t]
\centering
\includegraphics[width=\columnwidth]{figures/fig3_fpr_comparison.png}
\caption{False positive rate by prompting strategy and model on patched
contracts. Dashed line: FPR~$= 0.10$ acceptability threshold.}
\label{fig:fpr}
\end{figure}

\subsection{Tool Coverage Gap}

Figure~\ref{fig:coverage} places this work in context, showing the EVM-vs-non-EVM
coverage disparity that motivated the
benchmark~\cite{li2025solana}.

\begin{figure}[!t]
\centering
\includegraphics[width=\columnwidth]{figures/fig4_coverage_gap.png}
\caption{Number of published LLM-based audit tools by blockchain ecosystem.
Ethereum is saturated; Solana is partially covered; Algorand has effectively
no dedicated LLM tooling. Source:~\cite{li2025solana}.}
\label{fig:coverage}
\end{figure}

\subsection{Experimental Pipeline}

Figure~\ref{fig:pipeline} shows the end-to-end experimental pipeline.

\begin{figure}[!t]
\centering
\includegraphics[width=\columnwidth]{figures/fig5_pipeline_diagram.png}
\caption{Experimental pipeline: contracts and RAG corpus feed into prompting
strategies, which query three LLMs to produce 216 JSON outputs, scored to
produce the results table and figures.}
\label{fig:pipeline}
\end{figure}

\subsection{Explanation Quality Score Distribution}

Claude Sonnet~4 leads on explanation quality with a mean EQS of 4.24 across all
strategies, compared to 3.01 for Llama-3.3-70B and 2.88 for GPT-4o.
The CoT strategy yields the highest mean EQS of 3.62, ahead of 3.25 for both
zero-shot and RAG.
RC is 100\% under CoT for all models. Llama-3.3-70B shows the sharpest drop under
zero-shot, falling to 75\% RC, consistent with the highest EVM hallucination rate
on Algorand contracts.

Figure~\ref{fig:eqs} shows EQS distributions across prompting strategies for
correctly detected vulnerabilities only (DR~$= 1$ runs).

\begin{figure}[!t]
\centering
\includegraphics[width=\columnwidth]{figures/fig6_eqs_distribution.png}
\caption{Distribution of explanation quality scores (EQS 1 to 5) by prompting
strategy, restricted to correctly detected runs (DR~$= 1$).}
\label{fig:eqs}
\end{figure}

% -----------------------------------------------------------------------
\section{Discussion}
\label{sec:discussion}

\subsection{Prompting Strategy Is the Dominant Factor}

Across every metric, how the model is prompted matters more than which model
is used.
The gap between zero-shot and CoT is large, particularly for Algorand, where
pre-training coverage of PyTEAL is thin and models lack the scaffolding to
reason systematically through AVM-specific vulnerability classes.
Without explicit enumeration of what to look for, models tend toward superficial
pattern matching.
CoT resolves this by requiring the model to reason explicitly about each
vulnerability class in sequence before committing to a verdict.
RAG adds retrieved vulnerability documentation to the prompt, which helps most
on Algorand where pre-training gaps are largest, but at the cost of elevated
false positive rates: models shown patterns of what can go wrong appear more
likely to flag patched code as vulnerable.

There is a practical point here worth making directly.
When a developer pastes Solana or Algorand code into a general-purpose LLM
chat interface and asks for a security review, the interaction is functionally
a zero-shot query.
The detection rates in Section~\ref{sec:results} tell us what to expect from
that workflow.
Adding a CoT prompt that enumerates known vulnerability classes for the target
chain takes perhaps one extra paragraph of text and costs nothing in API terms.
The detection improvement is substantial.

\subsection{Chain Matters: Algorand Remains Harder}

All three models perform less well on Algorand than on Solana.
This is consistent with Boi and Esposito~\cite{boi2025prompt}, who found that
``fine-tuning offers benefits primarily when dealing with abstract, lower-level
languages such as TEAL (Algorand), where model comprehension is more limited
under prompt engineering alone.''
What this study adds is that the same difficulty applies to GPT-4o and
Claude Sonnet~4 under prompt engineering, not just open-source models.
Even the two strongest models in the evaluation show measurable difficulty on
V6 (LogicSig abuse) and V7 (group transaction manipulation).

The underlying issue is representation in pre-training corpora.
AVM execution semantics---atomic transaction groups, LogicSig statelessness,
flat fee model---appear far less frequently in the data these models were
trained on than Rust or Solidity semantics do.
RAG's largest relative benefit in this study is on Algorand, which makes sense:
it is injecting exactly the context that pre-training failed to provide.

\subsection{EVM Hallucinations Are a Real Problem}

The RC metric captures something that DR and FPR miss: model outputs that are
wrong in ways a practitioner might not immediately notice.
Across Algorand runs, examples included describing TEAL opcodes as
``Solidity-like,'' raising reentrancy as an Ethereum-style concern for AVM
contracts, and flagging gas-based denial-of-service risks in a system with a
flat fee model.
These errors are not benign.
A developer who receives a security report flagging non-existent runtime
properties may waste time investigating phantom risks, or worse, may be
distracted from an actual vulnerability that does not fit the EVM frame.
Llama-3.3-70B produces the highest hallucination rate on Algorand in zero-shot
conditions.
Under CoT, all three models achieve 100\% RC, suggesting that systematic
reasoning is also a partial corrective for cross-chain contamination.

\subsection{Performance Ceiling and Residual Hard Cases}

V4 (bump seed canonicalization) and V5 (stale CPI data) are the benchmark's
hardest cases.
The vulnerable code in both classes compiles and runs correctly; the flaw only
becomes apparent when reasoning about cross-instruction state or the semantics
of PDA derivation.
There is no syntactic signal a surface-level scanner could catch.
Even under CoT, these classes require reasoning about what the program
\emph{implies} about system state, not just what it says in the instruction
body.
Combining LLMs with lightweight static analysis---VRust-style dataflow analysis
for Solana is the natural candidate---seems necessary to reliably cover these
cases.

% -----------------------------------------------------------------------
\section{Threats to Validity}
\label{sec:threats}

\textbf{Internal validity.}
Both DR and RC scoring rely substantially on keyword matching.
This can produce false negatives when a model correctly identifies a
vulnerability but describes it in unexpected terms, and false positives when a
keyword is present without genuine detection.
The FPR heuristic is section-aware and negation-aware, which handles the common
cases, but edge cases remain.
Manual review of a random 20\% sample and inter-rater reliability checks for
EQS reduce the risk but do not eliminate it.

\textbf{Construct validity.}
EQS operationalizes explanation quality as the presence of four textual features:
vulnerability identification, code location, attack scenario, and a fix
recommendation.
This is a tractable proxy, not a complete measure of quality.
A model that gives a terse but technically correct fix could score lower than
a verbose one covering the same ground, and RC scoring does not catch all
forms of factual error.

\textbf{External validity.}
Twenty-four contract instances across eight vulnerability classes is a limited
evaluation surface.
The results may not hold for Solana and Algorand programs of different types,
for classes not included in the benchmark, or for production codebases at scale.
Each instance isolates a single vulnerability; real programs frequently contain
multiple interacting flaws that change how individual issues manifest.

\textbf{Conclusion validity.}
With 216 runs across 24 unique contract instances, statistical significance
testing is not meaningful.
All findings should be read as exploratory and direction-setting.
Confirmatory conclusions require replication at larger scale.

% -----------------------------------------------------------------------
\section{Conclusion}
\label{sec:conclusion}

This work provides, to our knowledge, the first systematic evaluation of frontier
LLMs on vulnerability detection for non-EVM smart contracts, covering Solana
and Algorand with a controlled dataset of 24 annotated contract pairs.
The results reveal three clear patterns.
Prompting strategy is a stronger determinant of detection accuracy than model
identity: CoT achieves 100\% DR and zero false positives regardless of which
model is used.
Algorand consistently proves harder than Solana for all three models, which we
attribute to limited PyTEAL representation in pre-training corpora.
And EVM-specific hallucinations---reentrancy claims, gas-cost concerns,
Solidity-like terminology---appear in a non-trivial fraction of Algorand runs
even for models that perform well on Solana.

For developers, the practical implication is straightforward: a CoT prompt that
enumerates the specific vulnerability classes of the target chain outperforms
generic zero-shot analysis substantially and costs nothing extra in API usage.
For the research community, the dataset, scoring scripts, and raw model outputs
released with this paper provide a reproducible starting point for further
work on non-EVM security analysis, an area that remains significantly
under-resourced relative to its Ethereum counterpart.

The benchmark has clear limitations---24 contracts, eight classes, isolated
single-vulnerability instances---and should not be taken as a comprehensive
evaluation. Future work ought to scale the contract corpus, add Move-based
chains such as Sui and Aptos, and investigate hybrid pipelines that pair LLM
reasoning with chain-specific static analysis to address the residual hard cases
that prompt engineering alone cannot reliably solve.

% -----------------------------------------------------------------------
\bibliographystyle{IEEEtran}
\bibliography{references}

\end{document}
